\chapter{BEBERAPA TEORI SPEKTRAL}\label{spectrum}

Dalam bab ini, untuk pertama kalinya kita akan menganggap, kecuali dinyatakan lain, bahwa ruang
vektor dan aljabar berskalar kompleks, bukan real. Sebagian besar alasan untuk hal ini
 \index{konvensi!dalam bab~\ref{spectrum} skalar bersifat kompleks}%
adalah perlunya menggunakan \emph{teorema Liouville}~\ref{C073131} untuk membuktikan
proposisi~\ref{C073134}.

\section{Spektrum}

\begin{defn} Suatu elemen $a$ dari aljabar beridentitas $A$ disebut
 \index{invertibel!kiri!dalam aljabar}%
\df{invertibel kiri} jika terdapat elemen $a_l$ dalam $A$ (disebut \df{invers kiri}
dari~$a$) sedemikian sehingga $a_l a = \vc 1$, dan disebut
 \index{invertibel!kanan!dalam aljabar}%
\df{invertibel kanan} jika terdapat elemen $a_r$ dalam $A$ (disebut \df{invers
kanan} dari~$a$) sedemikian sehingga $aa_r = \vc 1$. Elemen tersebut disebut
 \index{invertibel!dalam aljabar}%
\df{invertibel} jika elemen itu sekaligus invertibel kiri dan invertibel kanan. Himpunan semua
elemen invertibel dari $A$ dinotasikan
 \index{invertibel@$\inv A$ (elemen invertibel dalam aljabar)}%
dengan~$\inv A$.
\end{defn}

Suatu elemen dari aljabar beridentitas memiliki paling banyak satu invers perkalian. Bahkan,
berlaku pernyataan yang lebih kuat berikut.

\begin{prop}\label{000521} Jika suatu elemen dari aljabar beridentitas memiliki invers kiri dan
invers kanan, maka kedua invers tersebut sama (sehingga elemen itu invertibel).
\end{prop}

Jika suatu elemen $a$ dari aljabar beridentitas invertibel, inversnya yang unik dinotasikan
dengan~$a^{-1}$.

\begin{prop} Jika $a$ elemen invertibel dari suatu aljabar beridentitas, maka $a^{-1}$ juga
 \index{invers!dari suatu invers}%
invertibel dan
   \[ \bigl(a^{-1}\bigr)^{-1} = a\,. \]
\end{prop}

\begin{prop} Jika $a$ dan $b$ elemen-elemen invertibel dari suatu aljabar beridentitas, maka hasil kali
 \index{invers!dari suatu hasil kali}%
$ab$ juga invertibel dan
  \[ (ab)^{-1} = b^{-1}a^{-1}\,. \]
\end{prop}

\begin{prop}\label{000524} Jika $a$ dan $b$ elemen-elemen invertibel dari suatu aljabar beridentitas,
 \index{invers!selisih antara}%
maka
  \[ a^{-1} - b^{-1} = a^{-1}(b - a)b^{-1}\,. \]
\end{prop}

\begin{prop} Misalkan $a$ dan $b$ elemen-elemen dari suatu aljabar beridentitas. Jika $ab$ dan
$ba$ keduanya invertibel, maka $a$ dan $b$ juga invertibel.
\end{prop}

\begin{proof}[\emph{Petunjuk untuk bukti}] Gunakan proposisi~\ref{000521}. \ns \end{proof}

\begin{cor} Misalkan $a$ elemen dari suatu aljabar-$*\,$ beridentitas. Maka $a$ invertibel jika dan hanya jika $a^*a$ dan $aa^*$
invertibel; dalam hal ini,
   \[ a^{-1} = \bigl(a^*a\bigr)^{-1}a^* = a^*\big(aa^*\bigr)^{-1}\,. \]
\end{cor}

\begin{prop}\label{000526} Misalkan $a$ dan $b$ elemen-elemen dari suatu aljabar beridentitas. Maka
$\vc 1 - ab$ invertibel jika dan hanya jika $\vc 1 - ba$ invertibel.
\end{prop}

\begin{proof} Jika $\vc 1 - ab$ invertibel, maka $\vc 1 + b(\vc 1 - ab)^{-1}a$ adalah
invers dari $\vc 1 - ba$.
\end{proof}

\begin{exam} Contoh paling sederhana dari aljabar kompleks ialah keluarga $\C$ dari bilangan
kompleks. Aljabar ini beridentitas sekaligus komutatif.
\end{exam}

\begin{exam}\label{000532} Jika $X$ suatu ruang topologis tak kosong, maka keluarga $\fml C(X)$
dari semua fungsi kontinu bernilai kompleks pada $X$ merupakan suatu
 \index{C@$\fml C(X)$!sebagai aljabar beridentitas}%
aljabar di bawah operasi penjumlahan, perkalian skalar, dan perkalian titik demi titik yang biasa.
Aljabar ini beridentitas sekaligus komutatif.
\end{exam}

\begin{exam} Jika $X$ suatu ruang Hausdorff kompak lokal yang tidak kompak, maka (di bawah
operasi titik demi titik) keluarga $\fml C_0(X)$ dari semua fungsi kontinu bernilai kompleks
pada $X$ yang lenyap di tak hingga merupakan aljabar komutatif.
Namun, keluarga ini \emph{bukan} suatu
 \index{C@$\fml C_0(X)$!sebagai aljabar tanpa identitas}%
aljabar beridentitas.
\end{exam}

\begin{exam}\label{000533} Keluarga
 \index{Ma@$\mathbf M_n = \M n\C$!matriks bilangan kompleks berukuran $n \times n$}%
$\mathbf M_n$ dari matriks bilangan kompleks berukuran $n \times n$ merupakan
 \index{Ma@$\mathbf M_n = \M n\C$!sebagai aljabar beridentitas}%
aljabar beridentitas di bawah operasi penjumlahan, perkalian skalar, dan perkalian matriks yang biasa.
Aljabar ini tidak komutatif jika $n > 1$.
\end{exam}

\begin{exam}\label{000533a} Misalkan $A$ suatu aljabar. Jadikan keluarga
 \index{Ma@$\M nA$!matriks elemen aljabar berukuran $n \times n$}%
$\M n{A}$ dari matriks elemen $A$ berukuran $n \times n$ sebagai suatu
 \index{Ma@$\M nA$!sebagai aljabar}%
aljabar dengan menggunakan aturan operasi matriks yang sama dengan aturan untuk~$\mathbf M_n$.
Dengan demikian, $\mathbf M_n$ hanyalah $\M n\C$. Aljabar $\M nA$ beridentitas jika dan hanya jika $A$ beridentitas.
\end{exam}

\begin{exam}\label{000534} Jika $V$ suatu ruang linear bernorma, maka
 \index{B@$\ofml B(V)$!sebagai aljabar beridentitas}%
$\ofml B(V)$ merupakan aljabar beridentitas. Jika $\dim V~>~1$, aljabar tersebut tidak komutatif.
\end{exam}

\begin{defn}\label{000535} Misalkan $a$ elemen dari aljabar beridentitas~$A$. Maka
 \index{spektrum}%
\df{spektrum} dari $a$, yang dinotasikan dengan
 \index{sigma@$\sigma_A(a)$ atau $\sigma(a)$ (spektrum dari~$a$)}%
$\sigma_A(a)$ atau cukup $\sigma(a)$, adalah himpunan semua bilangan kompleks $\lambda$ sedemikian sehingga
$a - \lambda\vc 1$ tidak invertibel.
\end{defn}

\begin{exam} Jika $z$ suatu elemen dari aljabar bilangan kompleks $\C$, maka
 \index{spektrum!dari bilangan kompleks}%
$\sigma(z) = \{z\}$.
\end{exam}

\begin{exam}\label{0005712} Misalkan $X$ suatu ruang Hausdorff kompak. Jika $f$ suatu elemen
 \index{spektrum!dari fungsi kontinu}%
dari aljabar $\fml C(X)$ yang terdiri atas fungsi-fungsi kontinu bernilai kompleks pada $X$, maka
spektrum $f$ adalah jangkauannya.
\end{exam}

\begin{exam} Misalkan $X$ suatu ruang topologis sembarang. Jika $f$ suatu elemen dari
 \index{spektrum!dari fungsi kontinu terbatas}%
aljabar $\fml C_b(X)$ yang terdiri atas fungsi-fungsi kontinu terbatas bernilai kompleks pada $X$, maka
spektrum $f$ adalah tutupan jangkauannya.
\end{exam}

\begin{exam} Misalkan $S$ suatu ruang ukur positif dan $f \in \lfs{\infty}{S}$ suatu
(kelas ekuivalensi) fungsi yang terbatas secara esensial pada~$S$. Maka
 \index{spektrum!dari fungsi terbatas secara esensial}%
spektrum $f$ adalah jangkauan esensialnya.
\end{exam}

\begin{exam} Keluarga $\mathbf M_3$ dari matriks bilangan kompleks berukuran $3 \times 3$ merupakan aljabar
beridentitas di bawah operasi matriks yang biasa. Spektrum matriks $\begin{bmatrix} 5
& -6 & -6 \\ -1 & 4 & 2 \\ 3 & -6 & -4 \end{bmatrix}$ adalah $\{1,2\}$.
\end{exam}

\begin{prop} Misalkan $a$ suatu elemen dari aljabar beridentitas sedemikian sehingga $a^2 = \vc 1$.
 \index{spektrum!dari elemen yang kuadratnya $1$}%
Maka salah satu dari hal berikut berlaku:
 \begin{enumerate}[label=\rm{(\alph*)}]
  \item $a = \vc 1$, dan dalam hal ini $\sigma(a) = \{1\}$; atau
  \item $a = -\vc 1$, dan dalam hal ini $\sigma(a) = \{-1\}$; atau
  \item $\sigma(a) = \{-1,1\}$.
 \end{enumerate}
\end{prop}

\begin{proof}[\emph{Petunjuk untuk bukti}] Pada (c), untuk membuktikan $\sigma(a) \subseteq \{-1,1\}$,
tinjau $\D\frac1{1 - \lambda^2}(a + \lambda \mathbf 1).$  \ns
\end{proof}

\begin{prop} Ingat bahwa suatu elemen $a$ dari aljabar disebut
 \index{idempoten}%
 \index{spektrum!dari elemen idempoten}%
\df{idempoten} jika $a^2 = a$. Misalkan $a$ suatu elemen idempoten dari aljabar beridentitas. Maka
salah satu dari hal berikut berlaku:
 \begin{enumerate}[label=\rm{(\alph*)}]
  \item $a = \vc 1$, dan dalam hal ini $\sigma(a) = \{1\}$; atau
  \item $a = \vc 0$, dan dalam hal ini $\sigma(a) = \{0\}$; atau
  \item $\sigma(a) = \{0,1\}$.
 \end{enumerate}
\end{prop}

\begin{proof}[\emph{Petunjuk untuk bukti}] Pada (c), untuk membuktikan $\sigma(a) \subseteq \{0,1\}$,
tinjau $\D \frac1{\lambda - \lambda^2}\bigl(a + (\lambda - 1)\vc 1\bigr).$ \ns
\end{proof}

\begin{prop} Misalkan $a$ suatu elemen invertibel dari aljabar beridentitas. Maka suatu bilangan kompleks
tak nol $\lambda$ termasuk dalam spektrum $a$ jika dan hanya jika $1/\lambda$ termasuk dalam
spektrum~$a^{-1}$.
\end{prop}

Proposisi berikut merupakan akibat langsung dari proposisi~\ref{000526}.

\begin{prop}\label{000575} Jika $a$ dan $b$ elemen-elemen dari suatu aljabar beridentitas, maka,
kecuali mungkin untuk $0$, spektrum $ab$ dan spektrum $ba$ sama.
\end{prop}

\begin{defn}\label{C035854} Suatu pemetaan $f \colon A \sto B$ antara aljabar-aljabar Banach disebut
 \index{Banach!aljabar!homomorfisme}%
 \index{homomorfisme!aljabar Banach}%
\df{homomorfisme (aljabar Banach)} jika pemetaan tersebut sekaligus merupakan homomorfisme
aljabar dan pemetaan kontinu antara $A$ dan $B$. Kita menotasikan kategori
 \index{BALG@$\cat{BALG}$!kategori}%
 \index{kategori!$\cat{BALG}$ sebagai}%
aljabar Banach dan homomorfisme aljabar kontinu dengan $\cat{BALG}$.
\end{defn}

\begin{prop} Misalkan $A$ dan $B$ aljabar-aljabar Banach. Maka
 \index{langsung!jumlah!aljabar Banach}%
 \index{Banach!aljabar!jumlah langsung}%
\df{jumlah langsung} dari $A$ dan $B$, yang dinotasikan dengan $A \oplus B$, didefinisikan sebagai himpunan $A
\times B$ yang dilengkapi dengan operasi titik demi titik:
   \begin{enumerate}[label=\rm{(\alph*)}]
     \item $(a,b) + (a',b') = (a + a',b + b')$;
     \item $(a,b)\,(a',b') = (aa',bb')$;
     \item $\alpha(a,b) = (\alpha a,\alpha b)$
   \end{enumerate}
untuk $a$,$a' \in A$, $b$,$b' \in B$, dan $\alpha \in \C$. Seperti dalam~\ref{C023414}, definisikan
$\norm{(a,b)} = \norm a + \norm b$. Hal ini menjadikan $A \oplus B$ suatu aljabar Banach.
\end{prop}

\begin{exer} Identifikasi hasil kali dan koproduk (jika ada) dalam kategori $\cat{BALG}$
yang objeknya aljabar Banach
 \index{hasil kali!dalam $\cat{BALG}$}%
 \index{Banach!aljabar!hasil kali}%
 \index{Banach!aljabar!koproduk}%
 \index{koproduk!dalam $\cat{BALG}$}%
 \index{BALG@$\cat{BALG}$!hasil kali dan koproduk dalam}%
dan morfismenya homomorfisme aljabar kontinu, serta dalam kategori yang objeknya aljabar Banach dan
morfismenya homomorfisme aljabar kontraktif. (Di sini,
 \index{kontraktif}%
\emph{kontraktif} berarti $\norm{Tx} \le \norm x$ untuk setiap~$x$.)
\end{exer}

\begin{prop} Dalam suatu aljabar Banach $A$, operasi penjumlahan dan perkalian (yang dipandang sebagai
pemetaan dari $A \oplus A$ ke $A$) bersifat kontinu, dan perkalian skalar (yang dipandang sebagai pemetaan dari
$\K \oplus A$ ke $A$) juga bersifat kontinu.
\end{prop}

\begin{prop}[Deret Neumann]\label{prop_Neumann_series} Misalkan $a$ suatu elemen dari aljabar Banach beridentitas.
 \index{deret Neumann}%
 \index{deret!Neumann}%
Jika $\norm a < 1$, maka
$\vc 1 - a \in \inv A$ dan $(\vc 1 - a)^{-1} = \sum_{k=0}^\infty a^k$.
\end{prop}

\begin{proof}[\emph{Petunjuk untuk bukti}] Tunjukkan bahwa barisan $(\vc 1, a, a^2, \dots)$
dapat dijumlahkan. (Dalam aljabar beridentitas kita mengartikan $a^0$ sebagai~$\vc 1$.)  \ns
\end{proof}

\begin{cor}\label{cor_Neumann_series} Jika $a$ suatu elemen dari aljabar Banach beridentitas dengan $\norm{\vc 1 - a} < 1$, maka $a$ invertibel dan
    \[ \norm{a^{-1}} \le \frac{1}{1 - \norm{\vc 1 - a}}\,.  \]
\end{cor}

\begin{cor}\label{cor2_Neumann_series} Misalkan $a$ suatu elemen invertibel dari aljabar Banach beridentitas $A$ dan $b \in A$,
serta $\norm{a - b} < \norm{a^{-1}}^{-1}$. Maka $b$ invertibel dalam~$A$.
\end{cor}

\begin{proof}[\emph{Petunjuk untuk bukti}] Gunakan korolari sebelumnya untuk menunjukkan bahwa $a^{-1}b$ invertibel.
\ns \end{proof}

\begin{cor} Jika $a$ termasuk dalam suatu aljabar Banach beridentitas dan $\norm a < 1$, maka
   \[ \norm{(\vc 1 - a)^{-1} - \vc 1} \le \frac{\norm a}{1 - \norm a}\,. \]
\end{cor}

\begin{prop}\label{000644fa} Jika $A$ aljabar Banach beridentitas, maka $\open{\inv A}A$.
\end{prop}

\begin{proof}[\emph{Petunjuk untuk bukti}] Misalkan $a \in \inv A$. Tunjukkan, untuk $h$ yang cukup kecil,
bahwa $\vc 1 - a^{-1}h$ invertibel.  \ns
\end{proof}

\begin{prop}\label{000646fa} Misalkan $A$ aljabar Banach beridentitas. Pemetaan $a \mapsto a^{-1}$ dari $\inv A$ ke
dirinya sendiri merupakan homeomorfisme.
\end{prop}

\begin{prop} Misalkan $A$ aljabar Banach beridentitas. Pemetaan $r\colon a \mapsto a^{-1}$ dari
$\inv A$ ke dirinya sendiri bersifat diferensiabel dan pada setiap elemen invertibel $a$, berlaku
$dr_a(h) = -a^{-1}ha^{-1}$ untuk setiap $h \in A$.
\end{prop}

\begin{proof}[\emph{Petunjuk untuk bukti}] Untuk pengantar singkat mengenai diferensiasi pada ruang
berdimensi tak hingga, lihat Bab 13 dari catatan saya~\cite{Erdman:2007} tentang Analisis Real.  \ns
\end{proof}

\begin{prop}\label{000661} Misalkan $a$ suatu elemen dari aljabar Banach beridentitas~$A$. Maka spektrum $a$
kompak dan $\abs\lambda \le \norm a$ untuk setiap $\lambda \in \sigma(a)$.
\end{prop}

\begin{proof}[\emph{Petunjuk untuk bukti}] Gunakan \emph{teorema Heine-Borel}. Untuk membuktikan bahwa
spektrum tertutup, perhatikan bahwa $(\sigma(a))^c = f^\gets(\inv A)$ dengan $f(\lambda) = a -
\lambda\mathbf 1$ untuk setiap bilangan kompleks~$\lambda$. Tunjukkan juga bahwa jika $\abs\lambda >
\norm a$, maka $\mathbf 1 - \lambda^{-1}a$ invertibel.  \ns
\end{proof}

\begin{defn} Misalkan $a$ suatu elemen dari aljabar Banach beridentitas.
 \index{resolven!pemetaan}%
\df{Pemetaan resolven} untuk $a$ didefinisikan oleh
   \[ R_a \colon \C \setminus \sigma(a) \sto A
                 \colon \lambda \mapsto (a  - \lambda\vc 1)^{-1}\,. \]
\end{defn}

\begin{defn} Misalkan $\open U\C$ dan $A$ suatu aljabar Banach beridentitas. Suatu fungsi $f \colon U \sto A$
 \index{analitik}%
bersifat \df{analitik} pada $U$ jika
   \[ f'(a) := \lim_{z \sto a} \frac{f(z) - f(a)}{z - a} \]
ada untuk setiap $a \in U$. Suatu fungsi bernilai kompleks yang analitik pada seluruh $\C$
disebut fungsi
 \index{entire}
\df{entire}.
\end{defn}

\begin{prop} Untuk $a$ suatu elemen dari aljabar Banach beridentitas $A$ dan $\phi$ suatu fungsional linear
terbatas pada~$A$, misalkan $f := \phi \circ R_a \colon \C \setminus \sigma(a) \sto \C$. Maka
 \begin{enumerate}[label=\rm{(\alph*)}]
  \item $f$ analitik pada domainnya; dan
  \item $f(\lambda) \sto 0$ ketika $\abs\lambda \sto \infty$.
 \end{enumerate}
\end{prop}

\begin{proof}[\emph{Petunjuk untuk bukti}] Untuk (i), perhatikan bahwa dalam aljabar beridentitas $a^{-1} - b^{-1}
= a^{-1}(b - a)b^{-1}$.   \ns
\end{proof}

\begin{thm}[Liouville]\label{C073131} Setiap fungsi entire yang terbatas pada $\C$ bersifat konstan.
 \index{teorema Liouville}%
\end{thm}

\begin{proof}[\emph{Komentar tentang bukti}] Pembuktian teorema ini memerlukan sedikit latar belakang
dalam teori fungsi analitik, yakni materi yang tidak dibahas dalam catatan ini. Teorema
tersebut tentu dibahas dalam setiap mata kuliah pengantar tentang variabel kompleks. Sebagai contoh, bukti yang
baik dapat ditemukan dalam \cite{Rudin:1987}, teorema 10.23.   \ns
\end{proof}

\begin{prop}\label{C073134} Spektrum setiap elemen dari aljabar Banach beridentitas tidak kosong.
\end{prop}

\begin{proof}[\emph{Petunjuk untuk bukti}] Gunakan argumen kontradiksi. Gunakan \emph{teorema Liouville}~\ref{C073131}
untuk menunjukkan bahwa $\phi \circ R_a$ konstan bagi setiap fungsional linear terbatas $\phi$ pada~$A$. Kemudian gunakan
(salah satu versi) \emph{teorema Hahn-Banach} untuk membuktikan bahwa $R_a$ konstan. Mengapa konstanta ini harus~$0$?   \ns
\end{proof}

\noindent\emph{Catatan.} Penting untuk diingat bahwa kita hanya bekerja dengan aljabar
\emph{kompleks}. Hasil ini \emph{salah} untuk aljabar Banach real. Contoh
tandingan sederhana adalah matriks $\begin{bmatrix}0 & 1 \\ -1 & 0 \end{bmatrix}$ yang dipandang sebagai
elemen aljabar Banach (real) $M_2$ yang terdiri atas semua matriks $2 \times 2$ dengan entri
bilangan real.

 \vskip 1 pt

\begin{thm}[Gelfand-Mazur]\label{C073141} Jika $A$ aljabar Banach beridentitas yang setiap
 \index{teorema Gelfand-Mazur}%
elemen taknolnya invertibel, maka $A$ isomorfik secara isometrik dengan~$\C$.
\end{thm}

\begin{proof}[\emph{Petunjuk untuk bukti}] Misalkan $B = \{ \lambda\vc 1\colon \lambda \in \C\}$. Gunakan
proposisi sebelumnya~\ref{C073134} untuk menunjukkan bahwa $B = A$.    \ns
\end{proof}

Pernyataan berikut merupakan konsekuensi langsung dari \emph{teorema Gelfand-Mazur}
\ref{C073141} dan proposisi~\ref{C037821}.

\begin{cor}\label{C073147} Jika $I$ ideal maksimal dalam aljabar Banach komutatif beridentitas $A$,
maka $A/I$ isomorfik secara isometrik dengan~$\C$.
\end{cor}

\begin{prop}\label{0006703fa} Misalkan $a$ suatu elemen dari aljabar beridentitas. Maka
   \[ \sigma(a^n) = [\sigma(a)]^n \]
untuk setiap $n \in \N$. (Notasi $[\sigma(a)]^n$ berarti $\{\lambda^n\colon \lambda \in
\sigma(a)\}$.)
\end{prop}

\begin{defn}\label{C073154} Misalkan $a$ suatu elemen dari aljabar beridentitas.
 \index{spektral!radius}%
 \index{radius!spektral}%
\df{radius spektral} dari~$a$, yang dinotasikan dengan $\rho(a)$, didefinisikan sebagai
$\sup\{\abs\lambda\colon \lambda \in \sigma(a)\}$.
\end{defn}

\begin{prop}\label{0006703} Jika $a$ suatu elemen dari aljabar Banach beridentitas, maka $\rho(a) \le
\norm a$ dan $\rho(a^n) = \bigl(\rho(a)\bigr)^n$.
\end{prop}

\begin{prop} Misalkan $\phi\colon A \sto B$ homomorfisme yang mempertahankan identitas di antara aljabar-aljabar beridentitas. Jika $a$ elemen invertibel dari $A$,
maka $\phi(a)$ invertibel dalam $B$ dan inversnya adalah~$\phi(a^{-1})$.
\end{prop}

\begin{cor} Jika $\phi\colon A \sto B$ homomorfisme yang mempertahankan identitas di antara aljabar-aljabar beridentitas, maka
   \[ \sigma(\phi(a)) \subseteq \sigma(a) \]
untuk setiap $a \in A$.
\end{cor}

\begin{cor} Jika $\phi\colon A \sto B$ homomorfisme yang mempertahankan identitas di antara aljabar-aljabar beridentitas, maka
   \[ \rho(\phi(a)) \le \rho(a) \]
untuk setiap $a \in A$.
\end{cor}

\begin{thm}[Rumus radius spektral]\label{C073157} Jika $a$ suatu elemen dari aljabar Banach
 \index{spektral!radius!rumus untuk}%
beridentitas, maka
   \[ \rho(a) = \inf\bigl\{\norm{a^n}^{1/n}\colon n \in \N\bigr\}
                                = \lim_{n \sto \infty} \norm{a^n}^{1/n}\,. \]
\end{thm}

\begin{exer} Dengan menggunakan rumus untuk operator Volterra $V$ dalam contoh~\ref{000319} sebagai
 \index{operator Volterra!dan persamaan integral}%
 \index{operator!Volterra!dan persamaan integral}%
acuan, pada ruang Banach~$\fml C([0,1])$ definisikan operator integral melalui rumus $Vf(x)=\int_0^x f(t)\,dt$.
 \begin{enumerate}
  \item[(a)] Hitung spektrum operator $V$. (\emph{Petunjuk.} Tunjukkan bahwa $\norm{V^n}
\le (n!)^{-1}$ untuk setiap $n \in \N$ dan gunakan \emph{rumus radius spektral}.)
  \item[(b)] Apa implikasi (a) terhadap kemungkinan memperoleh fungsi kontinu $f$ yang
memenuhi persamaan integral
   \[ f(x) - \mu\int_0^xf(t)\,dt = h(x), \]
dengan $\mu$ suatu skalar dan $h$ suatu fungsi kontinu yang diberikan pada~$[0,1]$?
  \item[(c)] Gunakan gagasan dalam (b) dan uraian deret Neumann untuk $\vc 1 - \mu V$ (lihat
proposisi~\ref{prop_Neumann_series}) untuk menghitung secara eksplisit (dan dalam bentuk fungsi elementer)
suatu solusi untuk persamaan integral
 \[f(x) - {\textstyle \frac12}\int_0^xf(t)\,dt = e^x.\]
 \end{enumerate}
\end{exer}

\begin{prop} Misalkan $A$ aljabar Banach beridentitas dan $B$ subaljabar tertutup yang memuat~$\vc 1_A$.
Maka
 \begin{enumerate}[label=\rm{(\alph*)}]
   \item $\inv B$ terbuka sekaligus tertutup dalam $B \cap \inv A$;
   \item $\sigma_A(b) \subseteq \sigma_B(b)$ untuk setiap $b \in B$; dan
   \item jika $b \in B$ dan $\sigma_A(b)$ tidak memiliki lubang (yakni, jika komplemennya dalam $\C$
terhubung), maka $\sigma_A(b) =    \sigma_B(b)$.
 \end{enumerate}
\end{prop}

\begin{proof}[\emph{Petunjuk untuk bukti}] Tinjau fungsi
$f_b\colon (\sigma_A(b))^c \sto B \cap \inv A\colon \lambda \mapsto b - \lambda \vc 1$.   \ns
\end{proof}















\section{Spektrum Operator Ruang Hilbert}
\begin{prop} Misalkan $H$ ruang Hilbert dan $T \in \ofml B(H)$ suatu operator swaadjoin. Pernyataan berikut
ekuivalen:
 \begin{enumerate}
    \item[(i)] $\sigma(T) \subseteq [0,\infty)$;
    \item[(ii)] terdapat operator $S$ pada $H$ sedemikian sehingga $T = S^*S$; dan
    \item[(iii)] $T$ positif.
 \end{enumerate}
\end{prop}

\begin{notn} Jika $A$ adalah himpunan bilangan kompleks, kita
 \index{<unopurstar@$A^*$ (himpunan konjugat kompleks dari $A \subseteq \C$)}%
mendefinisikan $A^* := \{\conj\lambda \colon \lambda \in A\}$. Hal ini dilakukan untuk menghindari kerancuan dengan tutupan himpunan~$A$.
\end{notn}

\begin{prop} Jika $T$ suatu operator ruang Banach, maka $\sigma(T^*) = \sigma(T)$, sedangkan jika operator tersebut merupakan operator ruang Hilbert,
maka $\sigma(T^*) = (\sigma(T))^*$.
\end{prop}

Dalam ruang berdimensi hingga, hanya ada satu cara bagi suatu bilangan kompleks $\lambda$ untuk menjadi anggota spektrum suatu operator~$T$.
Alasannya adalah bahwa hanya ada satu cara suatu operator, dalam hal ini $T - \lambda I$, pada ruang berdimensi hingga dapat gagal
menjadi invertibel (lihat paragraf sebelum definisi~\ref{defn_similar_vs_ops}). Di sisi lain, ada beberapa cara bagi $\lambda$
untuk menjadi anggota spektrum suatu operator pada ruang berdimensi tak hingga---karena ada beberapa cara suatu operator dapat gagal
menjadi invertibel. (Lihat, misalnya, proposisi~\ref{prop_nasc_op_invert}.) Secara historis, hal ini telah melahirkan sejumlah usulan
taksonomi spektrum suatu operator.

\begin{defn} Misalkan $T$ suatu operator pada ruang Hilbert (atau Banach)~$H$. Suatu bilangan kompleks $\lambda$ termasuk dalam
 \index{resolven!himpunan}%
\df{himpunan resolven} dari $T$ jika $T - \lambda I$ invertibel. Himpunan resolven adalah komplemen dari \emph{spektrum}
$T$. Himpunan bilangan kompleks $\lambda$ yang membuat $T - \lambda I$ tidak injektif disebut
 \index{spektrum titik}%
 \index{spektrum!titik}%
\df{spektrum titik} dari $T$ dan dinotasikan
 \index{sigma@$\sigma_p(T)$ (spektrum titik dari~$T$)}%
dengan $\sigma_p(T)$. Anggota spektrum titik dari $T$ adalah
 \index{nilai eigen}%
\df{nilai eigen} dari $T$. Himpunan bilangan kompleks $\lambda$ yang membuat $T - \lambda I$ tidak terbatas jauh dari nol disebut
 \index{spektrum titik aproksimatif}%
 \index{spektrum!titik aproksimatif}%
\df{spektrum titik aproksimatif} dari $T$ dan dinotasikan
 \index{sigma@$\sigma_{ap}(T)$ (spektrum titik aproksimatif dari~$T$)}%
dengan $\sigma_{ap}(T)$. Himpunan semua bilangan kompleks $\lambda$ sedemikian sehingga tutupan jangkauan $T - \lambda I$ merupakan subruang
sejati dari $H$ disebut
 \index{spektrum kompresi}%
 \index{spektrum!kompresi}%
\df{spektrum kompresi} dari $T$ dan dinotasikan
 \index{sigma@$\sigma_c(T)$ (spektrum kompresi dari~$T$)}%
dengan $\sigma_c(T)$. Adapun
 \index{spektrum residual}%
 \index{spektrum!residual}%
\df{spektrum residual} dari $T$, yang kita notasikan
 \index{sigma@$\sigma_r(T)$ (spektrum residual dari~$T$)}%
dengan $\sigma_r(T)$, adalah $\sigma_c(T) \setminus \sigma_p(T)$.
\end{defn}

\begin{prop} Jika $T$ suatu operator ruang Hilbert, maka $\sigma_p(T) \subseteq \sigma_{ap}(T) \subseteq \sigma(T)$.
\end{prop}

\begin{prop} Jika $T$ suatu operator ruang Hilbert, maka $\sigma(T) = \sigma_{ap}(T) \cup \sigma_c(T)$.
\end{prop}

\begin{prop} Jika $T$ suatu operator ruang Hilbert, maka $\sigma_p(T^*) = (\sigma_c(T))^*$.
\end{prop}

\begin{cor} Jika $T$ suatu operator ruang Hilbert, maka $\sigma(T) = \sigma_{ap}(T) \cup (\sigma_{ap}(T^*))^*$.
\end{cor}

\begin{thm}[Teorema Pemetaan Spektral]\label{SMThm} Jika $T$ suatu operator ruang Hilbert dan $p$ suatu polinomial, maka
   \[ \sigma(p(T)) = p(\sigma(T)). \]
\end{thm}

Teorema sebelumnya berlaku secara lebih umum untuk setiap fungsi analitik yang didefinisikan pada lingkungan spektrum~$T$.
(Untuk pembuktiannya, lihat~\cite{Conway:1990}, bab VII, teorema 4.10.) Hasil tersebut juga berlaku untuk bagian-bagian spektrum.

\begin{prop} Jika $T$ suatu operator ruang Hilbert dan $p$ suatu polinomial, maka
   \begin{enumerate}[label=\rm{(\alph*)}]
     \item $\sigma_p(p(T)) = p(\sigma_p(T))$,
     \item $\sigma_{ap}(p(T)) = p(\sigma_{ap}(T))$, dan
     \item $\sigma_c(p(T)) = p(\sigma_c(T))$,
   \end{enumerate}
\end{prop}

\begin{prop}
    Jika $T$ operator invertibel pada ruang Hilbert, maka $\sigma(T^{-1}) = (\sigma(T))^{-1}$.
\end{prop}

\begin{defn}\label{defn_similar_Hsp_ops} Seperti halnya pada ruang vektor, dua operator $R$ dan $T$ pada ruang Hilbert (atau Banach) $H$ dikatakan
 \index{serupa!operator}%
\df{serupa} jika terdapat operator invertibel $S$ pada $H$ sedemikian sehingga $R = STS^{-1}$.
\end{defn}

\begin{prop} Jika $S$ dan $T$ merupakan operator yang serupa pada ruang Hilbert, maka $\sigma(S) = \sigma(T)$, $\sigma_p(S) = \sigma_p(T)$,
$\sigma_{ap}(S) = \sigma_{ap}(T)$, dan $\sigma_c(S) = \sigma_c(T)$.
\end{prop}

\begin{prop} Jika $T$ merupakan operator normal pada ruang Hilbert, maka $\sigma_c(T) = \sigma_p(T)$, sehingga $\sigma_r(T)~=~\emptyset$.
\end{prop}

\begin{prop} Jika $T$ merupakan operator pada ruang Hilbert, maka $\sigma_{ap}(T)$ tertutup.
\end{prop}

\begin{exam} Misalkan $D = \diag(a_1,a_2,a_3,\dots)$ merupakan operator diagonal pada ruang Hilbert dan misalkan
 \index{spektrum!bagian-bagian!operator diagonal}%
 \index{operator!diagonal!bagian-bagian spektrum}%
 \index{diagonal!operator!bagian-bagian spektrum}%
$A = \{a_k\colon k\in\N\}$. Maka
   \begin{enumerate}[label=(\alph*)]
     \item $\sigma_p(D) = \sigma_c(D) = A$,
     \item $\sigma_{ap}(D) = \clo A$, dan
     \item $\sigma_r(D) = \emptyset$.
   \end{enumerate}
\end{exam}

\begin{exam} Misalkan $S$ merupakan operator pergeseran unilateral pada $l_2$
 \index{spektrum!bagian-bagian!pergeseran unilateral}%
 \index{operator!pergeseran unilateral!bagian-bagian spektrum}%
 \index{pergeseran unilateral!bagian-bagian spektrum}%
(lihat contoh~\ref{C063526}). Maka
    \begin{enumerate}[label=(\alph*)]
      \item $\sigma(S) = D$,
      \item $\sigma_p(S) = \emptyset$,
      \item $\sigma_{ap}(S) = C$,
      \item $\sigma_c(S) = B$,
      \item $\sigma(S^*) = D$,
      \item $\sigma_p(S^*) = B$,
      \item $\sigma_{ap}(S^*) = D$, dan
      \item $\sigma_c(S^*) = \emptyset$.
    \end{enumerate}
di mana $B$ adalah cakram satuan terbuka, $D$ adalah cakram satuan tertutup, dan $C$ adalah lingkaran satuan pada bidang kompleks.
\end{exam}

\begin{exam} Misalkan $(S,\sfml A, \mu)$ merupakan ruang ukuran positif $\sigma$-hingga, $\phi \in L_\infty(S,\mu)$, dan $M_\phi$ merupakan
 \index{spektrum!bagian-bagian!operator perkalian}%
 \index{operator!perkalian!bagian-bagian spektrum}%
 \index{operator perkalian!bagian-bagian spektrum}%
operator perkalian yang bersesuaian pada ruang Hilbert $\lfs2{S,\mu}$ (lihat contoh~\ref{C063527}). Maka
spektrum dan spektrum titik aproksimatif dari $M_\phi$ adalah jangkauan esensial~$\phi$, sedangkan spektrum titik
dan spektrum kompresi dari $M_\phi$ adalah $\{\lambda \in \C\colon \mu(\phi^{\gets}(\{\lambda\})) > 0\}$.
\end{exam}






















\endinput
